\section{Introduction}

Unbounded lists in cycles, are a fundamental concept in computer science and
mathematics, often used to model repetitive structures or processes. They can
be thought of as infinite lists that repeat a finite sequence of elements.

\begin{equation*}
\begin{aligned}
L &= [x_0, x_1, x_2, \ldots, x_{n-1}]
  \mid x_n \in \mathbb{S}, L \in \mathbb{L} \\
\operatorname{Cycle}(L)
  &= [x_0, x_1, x_2, \ldots, x_{n-1}, x_0, x_1, \ldots]
\end{aligned}
\end{equation*}

In this article, we present discrete definition of Cycle operations over finite
integer lists, defined recursively and verified some of its properties using
the Stainless system. Our approach follows a zero-prior-knowledge philosophy,
building on a previously verified foundation for recursive list structure. The
result is a verified, from-scratch implementation of cycle operations, suitable
as a foundation for higher-level numeric reasoning over unbounded lists.

This article verifies:

\begin{itemize}
  \item Cycle definitions: recursive, modulo, and memory --- Section~3.
  \item Equivalence: recursive and modulo produce identical values at every
    position --- Section~4.
  \item Element access: modular indexing, small-position direct lookup ---
    Subsections~5.1--5.2.
  \item Periodic invariance: value unchanged by adding cycle-period multiples
    --- Subsections~5.3--5.4.
  \item Mod propagation: remainder computed from base-cycle values ---
    Subsection~5.5.
  \item Repeated-cycle invariance: repeating the base list preserves all
    lookups --- Subsection~5.6.
  \item Cycle value positivity: all values $\geq 0$ at every position ---
    Subsection~5.7.
  \item Cycle rotation: rotates base list, shifts index --- Subsection~5.8.
  \item Residue classification: all-zero, none-zero, and some-zero divisor
    classifications transfer from the base list to every cycle position ---
    Subsections~5.10--5.12.
\end{itemize}

\subsection{Related work}

Finite rotation already has a substantial formal treatment in Lean's Mathlib:
it includes rotation reduced modulo list length, indexed rotation laws, and an
equivalence relation identifying rotated lists~\cite{leanmathlibrotate}. These
results give a close formal point of contact for the article's base-list
rotation and modular-indexing laws.

There is also broader work on cyclic datatypes as recursive programming
structures. Hamana studies cyclic lists and related data modulo bisimulation,
with a semantic account of their computation rules~\cite{hamana2017cyclic}.
That setting is more general and different from the present periodic sequence
generated by a finite integer list. Together, these sources place the Stainless
development between finite-list rotation and general cyclic-data theory, while
this article verifies the equivalence of its recursive and modulo presentations
and the period, residue, and lookup properties of its concrete periodic model.
